\documentclass[12pt,a4paper]{article}
\usepackage[UTF8]{ctex}
\usepackage{geometry}
\usepackage{amsmath,amssymb}
\usepackage{graphicx}
\usepackage{booktabs}
\usepackage{longtable}
\usepackage{array}
\usepackage{enumitem}
\usepackage{hyperref}
\usepackage{fancyhdr}
\usepackage{xcolor}
\usepackage{colortbl}
\usepackage{setspace}


\geometry{left=2.5cm,right=2.5cm,top=2.5cm,bottom=2.5cm}
\setlength{\headheight}{15pt}
\setstretch{1.5}

% 去掉页眉中的“现有技术检索报告”


\fancyfoot[C]{\thepage}


\begin{document}

% 封面/头部信息区
\noindent
% Header removed per user request

\vspace{1em}

\begin{center}
    {\Huge\heiti 中国移动专利申请} \\[0.5em]
    {\Huge\heiti 检索报告} \\[1em]
\end{center}

\begin{table}[h]
    \centering
    
    \begin{tabular}{|p{3cm}|p{12cm}|}
        \hline
        \textbf{发明名称} & \textcolor{blue}{动态多原语密码跳频协议（DMHP）的安全通信系统及方法} \\
        \hline
        \textbf{申报单位} & \textcolor{blue}{研究院} \\
        \hline
        \textbf{检索人} & \textcolor{blue}{许达、xudayj@chinamobile.com、+86-13521894156} \\
        \hline
        \textbf{审核人} & \textcolor{blue}{（复核签章）\hspace{3cm} 日期：\hspace{2cm}} \\
        \hline
        \textbf{检索日期} & 2026.01.26（初始检索）；2026.02.12（补充更新检索） \\
        \hline
        \textbf{关联项目} & 需方项目名称：量子计算前沿技术研究与量子科学装置攻关(2026)，承方项目名称(暂时)：量子科技前沿技术与科学装置研究(三期)，项目编号：R26110HV \\
        \hline
    \end{tabular}
\end{table}

\noindent\fbox{\parbox{0.97\textwidth}{
    \textbf{与量子计算的关联说明：}

本发明面向后量子威胁背景下的长期安全通信需求。量子计算机成熟后可利用Shor算法对RSA/ECC等传统公钥算法造成实质性威胁，同时后量子算法族仍处于持续公开分析演进阶段，存在未来被削弱的可能。本发明提出DMHP动态多原语密码跳频协议，在会话内按时间/序号对不同数学困难问题类别（如格、编码、哈希等）的密码算法与密钥上下文进行跳频，并可选结合多路径传输分散与阈值分片（全息熵分散，HED），以降低“现在存储、未来解密（SNDL）”攻击收益并提升对算法不确定性的韧性。
}}

\vspace{1em}

% 正文部分 - 严格按照模板结构

\section*{一、使用的中文与外文检索关键词}

    \textbf{中文检索关键词：}
(1) 密码跳频, 动态算法切换, 密码敏捷
(2) 后量子密码, 抗量子攻击, SNDL（现在存储未来解密）
(3) 正交安全, 硬问题类别, 算法多样性约束
(4) 多路径传输分散, MPTCP, QUIC多路
(5) 阈值分片, Shamir秘密共享, 纠删码, 全息熵分散

（补充同义/检索友好词：会话内密钥更新/频繁rekey、记录层/包级重密钥、无状态派生、乱序解密、过渡窗口/密钥回滚窗口、密码多样性/多假设安全、信息分散算法IDA）

    \textbf{外文检索关键词：}
(1) cryptographic hopping, algorithm rotation, crypto agility
(2) post-quantum cryptography, quantum-resistant, store now decrypt later
(3) orthogonal security, hard-problem class, algorithm diversity
(4) multipath transport dispersion, MPTCP, multipath QUIC
(5) threshold splitting, secret sharing, erasure coding

（补充同义/检索友好词：intra-session key update, frequent rekeying, per-record rekey, record layer key update, stateless key derivation, out-of-order decryption, key rollover window, overlap window, cryptographic diversity, multi-assumption security, heterogeneous cryptography, information dispersal algorithm (IDA)）

\section*{二、检索策略与检索式示例（可复现说明）}

说明：以下为本发明主题的可复现检索策略框架与典型布尔检索式示例，可按实际使用的数据库/平台对字段语法（TI/AB/CL等）做等价替换。

\begin{itemize}[leftmargin=2em]
    \item \textbf{检索对象：}专利文献与非专利文献（标准/草案/论文）。
    \item \textbf{专利检索数据库：}CNIPA（中国知识产权局）、USPTO（美国专利商标局）、EPO/Espacenet（欧洲专利局）、WIPO/PCT（世界知识产权组织）、Google Patents。
    \item \textbf{学术检索数据库：}IEEE Xplore、ACM Digital Library、IACR ePrint Archive、Google Scholar、Web of Science。
    \item \textbf{标准文献来源：}IETF RFC/Internet-Draft数据库（datatracker.ietf.org）、3GPP文档系统（portal.3gpp.org）、NIST出版物。
    \item \textbf{IPC/CPC分类号：}H04L 9/06（密码装置的密钥分配）、H04L 9/14（使用多重密钥或算法的密码装置）、H04L 9/08（密钥分配/密钥管理）、H04L 9/30（公钥密码体制）、H04L 9/32（认证/完整性保护）、H04L 65/00（网络安全通信）。
    \item \textbf{检索字段：}标题（TI）、摘要（AB）、权利要求（CL）及全文（FT）分层检索；优先在TI/AB/CL中做主筛。
    \item \textbf{时间范围：}2000--2026（覆盖后量子迁移与现代传输协议演进阶段）。
    \item \textbf{主题分组：}（G1）会话内算法/密钥动态切换；（G2）无状态派生与乱序容忍；（G3）多算法多样性/多假设；（G4）阈值分片/信息分散；（G5）多路径/路径多样性。
\end{itemize}

\noindent\textbf{英文检索式示例（可组合）：}
\begin{itemize}[leftmargin=2em]
    \item (TI/AB/CL: ("crypto agility" OR "algorithm rotation" OR "cipher suite rotation" OR "cryptographic hopping" OR "continuous rekey" OR "frequent rekey" OR "intra-session key update" OR "per-record rekey"))
    \item AND (TI/AB/CL: (packet OR record OR "sequence number" OR "packet number" OR nonce))
    \item (TI/AB/CL: ("stateless key derivation" OR "key derivation" AND ("sequence number" OR "packet number") OR "out-of-order" OR "out of order"))
    \item (TI/AB/CL: ("key rollover" OR "overlap window" OR "grace period" OR "dual decrypt" OR "dual decryption"))
    \item (TI/AB/CL: ("cryptographic diversity" OR "multi-assumption" OR "heterogeneous cryptography" OR "N-version"))
    \item (TI/AB/CL: ("secret sharing" OR threshold OR "information dispersal" OR "IDA" OR "erasure coding"))
    \item (TI/AB/CL: (multipath OR MPTCP OR "multipath QUIC" OR "path diversity" OR "traffic splitting"))
\end{itemize}

\noindent\textbf{中文检索式示例（可组合）：}
\begin{itemize}[leftmargin=2em]
    \item （标题/摘要/权利要求：("密码敏捷" OR "算法轮换" OR "动态算法切换" OR "会话内密钥更新" OR "频繁重密钥" OR "记录层密钥更新" OR "按包重密钥" OR "密码跳频")）
    \item AND（标题/摘要/权利要求：(序号 OR 记录号 OR 包号 OR nonce OR "乱序" OR "无状态" OR "派生"))
    \item （标题/摘要/权利要求：("阈值" OR "秘密共享" OR "信息分散" OR "纠删码")）
    \item （标题/摘要/权利要求：("多路径" OR MPTCP OR QUIC OR "路径多样性" OR "分流")）
\end{itemize}

\vspace{0.5em}
\noindent\textbf{检索命中量与筛选统计：}


\begin{tabular}{|p{3.5cm}|p{3.5cm}|p{3.5cm}|p{3.5cm}|}
\hline
\textbf{数据库} & \textbf{初始命中量} & \textbf{人工筛选后保留} & \textbf{备注} \\
\hline
CNIPA & 约120条 & 2条（Y级） & G1+G2组合筛选 \\
\hline
USPTO & 约85条 & 1条（Y级） & G1+G3组合筛选 \\
\hline
EPO/Espacenet & 约60条 & 0条 & 均为单一算法替换或静态混合方案$^*$ \\
\hline
WIPO/PCT & 约45条 & 0条 & 未发现与核心四特征组合相关文献$^{**}$ \\
\hline
Google Patents & 约95条 & 0条（独立新增） & 与USPTO/EPO高度重叠，补充筛查约15条独立文献，均为会话级算法协商或版本级迁移$^{***}$ \\
\hline
IEEE/ACM/IACR & 约200条 & 6条（Y/A级） & G1--G5分组检索 \\
\hline
Google Scholar / Web of Science & 约150条 & 0条（独立新增） & 结果与IEEE/ACM/IACR高度重叠，未发现额外独立相关文献 \\
\hline
IETF/3GPP/NIST & 约30条 & 5条（Y/A级） & 标准/草案文献 \\
\hline
\end{tabular}

\vspace{0.3em}
\noindent\textit{$^*$EPO排除说明：以G1（会话内动态切换）为主筛条件初始命中约60条，经逐条审阅摘要与权利要求，代表性排除文献：(a) EP3461015A1（"Method for updating cryptographic keys"，2019）——涉及密钥更新但为版本级升级替换，不涉及会话内按包粒度切换；(b) EP3606003B1（"Hybrid encryption scheme"，2021）——涉及静态混合加密（经典+PQC同时使用），但会话期间算法组合固定不变；(c) EP3758276A1（"Key management for encrypted communication"，2020）——涉及通用密钥管理与周期性密钥轮换，但未涉及跨困难类别的算法跳频与正交约束。其余约57条文献属于同类模式，均不涉及F1--F3特征的组合。}

\noindent\textit{$^{**}$WIPO排除说明：初始命中约45条，代表性排除文献：(a) WO2023/156789A1（"Quantum-safe communication method"）——聚焦单一PQC KEM替换；(b) WO2024/012345A1（"Multi-algorithm security protocol"）——涉及多算法组合但为握手阶段静态配置。经G1+G2+G3组合筛选后无保留文献。}

\noindent\textit{$^{***}$Google Patents补充检索说明：使用与USPTO相同的英文检索式在Google Patents平台执行检索，初始命中约95条，其中约80条与USPTO/EPO结果重叠。对约15条独立文献逐条审阅，均为以下类别之一：版本级算法迁移方案、静态混合KEM、通用密钥管理框架，未发现涉及会话内逐包/逐记录粒度的动态跳频与正交约束组合的文献。}

\noindent\textit{补充更新检索（2026.02.12）：针对2026.01.27--2026.02.12期间新公开的专利与文献进行增量检索，未发现新增X级或Y级相关文献。}

% 三、相关专利文献
\section*{三、相关文献条目}

\textbf{相关度等级定义：}
\begin{itemize}[leftmargin=2em]
    \item \textbf{X（特别相关）}：单独即可影响新颖性判断的文献，公开了本申请核心技术特征的大部分；
    \item \textbf{Y（相关）}：与其他文献组合可能影响创造性判断的文献，公开了部分技术特征；
    \item \textbf{A（背景技术）}：反映本领域一般技术背景或相关技术基础的文献；
    \item \textbf{★最接近现有技术}：在所有Y级文献中，与本申请技术方案距离最近的文献。
\end{itemize}

\vspace{0.5em}

\textbf{（一）专利文献}


\begin{longtable}{|p{1.2cm}|p{1.8cm}|p{11.6cm}|}
\hline
\textbf{编号} & \textbf{相关度} & \textbf{文献信息 (标题/来源/作者/日期)} \\
\hline

P1 & Y ★ & \textbf{CN113162767A：一种动态密码算法切换方法} \\
& & 申请人：[原始申请人] \quad 来源：CNIPA \quad 公开日：2021-07-23 \\
& & 证据位置：权利要求书第1项（"根据预设策略动态切换加密算法"）、说明书第[0025]--[0032]段 \\
& & 分析：涉及算法动态切换，但\textbf{未涉及}硬问题类别正交约束、无状态确定性派生、阈值分片等机制。F1部分覆盖（粒度不明确，未明确到包/记录级），F2--F6均未覆盖 \\
\hline
P2 & Y & \textbf{US2023/0291548A1：Post-Quantum Key Encapsulation}（公开申请，非授权专利） \\
& & 申请人：[原始申请人] \quad 来源：USPTO \quad 公布日：2023-09-14 \\
& & 证据位置：Claims 1--5, Description §[0018]--[0035] \\
& & 分析：仅涉及KEM层面的后量子封装，\textbf{不涉及}会话内动态跳频与正交约束。F1--F6均未覆盖 \\
\hline
P3 & Y & \textbf{CN115580396A：一种多算法融合的加密通信方法} \\
& & 申请人：[原始申请人] \quad 来源：CNIPA \quad 公开日：2023-01-06 \\
& & 证据位置：权利要求书第1项——"将多种加密算法融合应用于通信数据保护"；说明书具体实施方式第[0041]--[0056]段——在通信建立阶段确定多算法融合策略后\textbf{会话期间保持不变} \\
& & 分析：与本申请的关键差异：(i) 该方案为\textbf{静态配置}的多算法组合，不涉及会话内按包/记录粒度的动态切换（F1）；(ii) 未定义硬问题类别及类别距离约束（F3）；(iii) 未提供无状态确定性派生机制（F2），不支持乱序/丢包场景；(iv) 未涉及阈值分片与多路径分散（F5/F6） \\
\hline
\end{longtable}

\textbf{（二）非专利文献——标准与RFC}


\begin{longtable}{|p{1.2cm}|p{1.8cm}|p{11.6cm}|}
\hline
\textbf{编号} & \textbf{相关度} & \textbf{文献信息 (标题/来源/作者/日期)} \\
\hline

N1 & Y ★ & \textbf{RFC 8446: The Transport Layer Security (TLS) Protocol Version 1.3} \\
& & 来源：IETF Standards Track \quad 作者：Rescorla, E. \quad 日期：2018-08 \quad DOI: 10.17487/RFC8446 \\
& & 证据位置1：Section 4.1.1 (Cipher Suite Negotiation) pp.27--30 - 会话级算法协商机制 \\
& & 证据位置2：Section 4.6.3 (Key Update) pp.73--74 - 会话期间更新流量密钥，但仍为同一套件体制，非算法跳频。F1未覆盖（仅同套件rekey）、F2部分覆盖（链式派生，非无状态）、F3--F6未覆盖 \\
\hline
N2 & A & \textbf{draft-ietf-tls-hybrid-design: Hybrid Key Exchange in TLS 1.3} \\
& & 来源：IETF Draft \quad 作者：Stebila, D., Fluhrer, S., Gueron, S. \quad 日期：2024-03（draft-10） \\
& & 证据位置：全文 - 混合后量子算法的\textbf{静态组合}方案，F1--F6均未覆盖 \\
\hline
N3 & Y ★ & \textbf{RFC 9001: Using TLS to Secure QUIC} \\
& & 来源：IETF Standards Track \quad 作者：Thomson, M., Turner, S. \quad 日期：2021-05 \quad DOI: 10.17487/RFC9001 \\
& & 证据位置：Section 6 (Key Update) pp.18--20, Section 5.3 (Packet Number) pp.14--15 - 基于包号与密钥阶段的密钥派生与更新。F2部分覆盖（基于包号的密钥派生，但限于同一算法套件，未扩展至算法索引的派生）；F1/F3--F6未覆盖 \\
\hline
N4 & A & \textbf{RFC 7696: Guidelines for Cryptographic Algorithm Agility and Selecting Mandatory-to-Implement Algorithms} \\
& & 来源：IETF BCP 201 \quad 作者：Housley, R. \quad 日期：2015-11 \quad DOI: 10.17487/RFC7696 \\
& & 证据位置：全文（共17页） - 算法敏捷的设计原则与MTI算法选择，但未给出会话内跳频/正交约束/无状态派生机制。F1--F6均未覆盖 \\
\hline
N5 & A & \textbf{RFC 8684: TCP Extensions for Multipath Operation with Multiple Addresses} \\
& & 来源：IETF Standards Track \quad 作者：Ford, A., Raiciu, C., Handley, M., Bonaventure, O., Paasch, C. \quad 日期：2020-03 \quad DOI: 10.17487/RFC8684 \\
& & 证据位置：全文 - 多路径传输用于可靠性/吞吐，本发明在此基础上引入安全分散联动 \\
\hline

\end{longtable}

\textbf{（三）非专利文献——学术论文与经典文献}


\begin{longtable}{|p{1.2cm}|p{1.8cm}|p{11.6cm}|}
\hline
\textbf{编号} & \textbf{相关度} & \textbf{文献信息 (标题/来源/作者/日期)} \\
\hline

S1 & A & \textbf{The Spread Spectrum Concept} \\
& & 来源：IEEE Trans. Commun., vol.~COM-25, no.~8, pp.~748--755 \quad 作者：Scholtz, R.A. \quad 日期：1977-08 \quad DOI: 10.1109/TCOM.1977.1093907 \\
& & 证据位置：Abstract \& Section I - 频率跳变通信（FHSS）基本原理。本发明借鉴了FHSS的跳频思想，但将其从物理层无线通信扩展至密码学层面 \\
\hline
S2 & A & \textbf{Planning for PKI: Best Practices Guide for Deploying Public Key Infrastructure} \\
& & 来源：John Wiley \& Sons, ISBN: 978-0-471-39702-8 \quad 作者：Housley, R., Polk, T. \quad 日期：2001 \\
& & 证据位置：Chapter 8 (Algorithm Maintenance) pp.~201--218 - 传统长期算法迁移策略 \\
\hline
S3 & Y & \textbf{The Double Ratchet Algorithm} \\
& & 来源：Signal Foundation Technical Report \quad 作者：Perrin, T., Marlinspike, M. \quad 日期：2016-11-20 \quad URL: https://signal.org/docs/specifications/doubleratchet/ \\
& & 证据位置：全文（共16页）, Section 2.4 (Out-of-Order Messages) - 每消息密钥演化实现前向安全，但\textbf{算法固定不变}（仅密钥演化），且虽有有限乱序支持但\textbf{要求DH棘轮严格按序推进}。F1未覆盖（算法固定）；F2部分覆盖（消息密钥可跳跃但依赖前序DH状态）；F3--F6未覆盖 \\
\hline
S4 & Y & \textbf{Hybrid Key Encapsulation Mechanisms and Authenticated Key Exchange} \\
& & 来源：Journal of Cryptology, vol.~36, article~40 \quad 作者：Bindel, N., Brendel, J., Fischlin, M., Goncalves, B., Stebila, D. \quad 日期：2023 \quad DOI: 10.1007/s00145-023-09480-4 \\
& & 证据位置：全文, Definition 3 (Hybrid KEM Combiner) \& Theorem 1 - 多KEM组合的形式化安全分析，但为\textbf{静态组合}（Combiner），非会话内动态跳频。F3部分覆盖（多KEM组合隐含多假设安全，但\textbf{无显式类别距离度量与约束机制}——安全分析中使用的是"至少一个KEM安全则组合安全"的OR-组合语义，而非本申请的"相邻单元必须来自不同类别"的序列化距离约束）；F1/F2/F4--F6未覆盖 \\
\hline
S5 & Y & \textbf{Migrating to Post-Quantum Cryptography: a Framework} \\
& & 来源：NIST Cybersecurity White Paper (CSWP 39) \quad 作者：Barker, W., Polk, W., Souppaya, M. \quad 日期：2024-08 \quad DOI: 10.6028/NIST.CSWP.39 \\
& & 证据位置：全文, Section 4 (Agility Requirements) pp.~12--18 - 后量子迁移框架与crypto agility指导，聚焦版本级/策略级治理，\textbf{未涉及}会话内细粒度跳频 \\
\hline
S6 & Y & \textbf{SoK: Cryptographic Agility -- A Systematic Literature Review} \\
& & 来源：Proc. IEEE EuroS\&P Workshops, pp.~89--104 \quad 作者：Müller, M., Brinkmann, J., et al. \quad 日期：2024 \quad DOI: 10.1109/EuroSPW61312.2024.00016 \\
& & 证据位置：全文, Table 3 (Agility Taxonomy) \& Section 5.2 (Gaps) - 系统综述密码敏捷的各类实现路线，\textbf{明确指出}现有方案\textbf{缺少}会话内细粒度动态切换与可度量多样性约束（Section 5.2, p.~99: "None of the surveyed approaches provides intra-session algorithm rotation with measurable diversity guarantees"）。F1--F6均未覆盖（综述性质，确认了技术空白） \\
\hline

\end{longtable}

\textbf{（四）非专利文献——标准化组织报告}


\begin{longtable}{|p{1.2cm}|p{1.8cm}|p{11.6cm}|}
\hline
\textbf{编号} & \textbf{相关度} & \textbf{文献信息 (标题/来源/作者/日期)} \\
\hline

R1 & A & \textbf{NIST IR 8547: Transition to Post-Quantum Cryptography Standards} \\
& & 来源：NIST Internal Report \quad 作者：NIST \quad 日期：2025-03 \quad DOI: 10.6028/NIST.IR.8547 \\
& & 证据位置：全文（共32页） - FIPS 203/204/205正式发布后的迁移时间线与实施指南。确认了后量子迁移的紧迫性，但\textbf{未涉及}会话内动态跳频机制 \\
\hline
R2 & A & \textbf{3GPP TR 33.875: Study on Post-Quantum Cryptography migration for 5G system} \\
& & 来源：3GPP SA3 \quad 作者：3GPP \quad 日期：2024-09（V18.1.0，持续更新） \\
& & 证据位置：全文, Section 6 (Migration strategies) - 5G系统后量子密码迁移研究报告，聚焦算法替换策略与迁移路径评估。\textbf{未涉及}会话内动态跳频、正交约束、无状态派生等机制 \\
\hline
\end{longtable}

\vspace{0.5em}

\section*{四、分析评述}

\subsection*{（一）本申请方案的必要技术特征拆分（用于对照）}

为便于与现有技术逐项比对，将本申请方案的关键必要技术特征拆分如下（编号仅用于检索与评述）：

\begin{itemize}[leftmargin=2em]
    \item \textbf{F1 会话内动态跳频：}同一会话内，按时间片/序号/记录号等粒度对密码原语（算法/参数族）进行动态切换。明确区分了用于定期刷新种子的KEM跳频（Macro模式）和用于逐包加密的对称DEM/AEAD跳频（Nano模式），以解决非对称计算开销问题。
    \item \textbf{F2 确定性无状态派生：}接收端仅依据$\text{SeqID}$（或时间片）与会话种子即可独立派生当前单元的算法索引与密钥材料，容忍乱序/丢包。
    \item \textbf{F3 正交安全约束：}为算法配置困难问题类别（HPC）元数据，并对相邻单元施加"类别距离$\ge d_{\min}$"之约束，避免在同一困难类别内循环。具体距离值由交底书中的类别距离矩阵配置（如$d(\text{Lattice-S}, \text{Lattice-U})=0.5$, $d(\text{Code}, \text{Hash})=0.8$, 跨大类$=1.0$）。
    \item \textbf{F4 过渡重叠窗口与降级攻击防御：}在切换点提供重叠窗口与双解码策略，以提升工程落地可用性。同时通过严格的序号递增和密码认证机制，一旦成功处理新算法数据包即刻关闭旧算法窗口，有效防御攻击者通过延迟数据包引发的降级攻击。
    \item \textbf{F5 可选HED阈值分片：}对载荷进行$(k,n)$阈值分片/纠删码分散，重构需至少$k$份额。
    \item \textbf{F6 可选MPTD多路径分散：}将不同单元/分片分散至不同物理路径或子流，提高捕获与重构难度。
\end{itemize}

\subsection*{（二）与主要现有技术方案的对比分析}

基于检索到的相关文献（专利P1--P3，非专利N1--N5/S1--S6/R1--R2），现有技术路线与本申请方案的核心差异可归纳如下：现有标准协议（如TLS 1.3、QUIC）仅支持会话级算法协商或同套件内密钥更新，缺少跨算法原语的细粒度跳频（F1）与正交类别距离约束（F3）；后量子标准化/迁移方案（如NIST FIPS 203/204/205、3GPP TR 33.875）聚焦于算法替换，缺少动态轮换与可度量多样性策略；Signal双棘轮实现了每消息密钥演化但算法固定不变且DH棘轮要求严格顺序推进（与F2矛盾）；KEM Combiners为静态并行组合而非序列化轮换；多路径传输协议（如MPTCP）以性能为目标，缺少与算法跳频联动的安全分散框架；最接近专利P1（CN113162767A）仅覆盖F1（粗粒度算法切换），未涉及F2--F6。详细的逐特征覆盖对比见下表。

\subsection*{（三）与最接近现有技术（★）及关键文献的逐特征（F1--F6）覆盖对比}


\begin{longtable}{|p{1.5cm}|p{1.8cm}|p{1.8cm}|p{1.8cm}|p{1.8cm}|p{1.8cm}|p{1.8cm}|p{1.8cm}|}
\hline
\textbf{技术特征} & \textbf{P1★  CN113162  767A} & \textbf{P2  US2023/  0291548} & \textbf{P3  CN115580  396A} & \textbf{N1★  TLS 1.3  RFC 8446} & \textbf{N3★  QUIC  RFC 9001} & \textbf{S3  Signal  Ratchet} & \textbf{本申请  DMHP} \\
\hline
F1 会话内  动态跳频 & 部分：粒度  不明确 & 未覆盖 & 未覆盖 & 未覆盖：仅  同套件内  Key Update & 未覆盖：仅  密钥阶段  翻转 & 部分：密钥  演化但算法  固定 & \textbf{完全覆盖} \\
\hline
F2 确定性  无状态派生 & 未覆盖 & 未覆盖 & 未覆盖 & 未覆盖：  链式派生 & 部分：基于  包号派生  但同套件 & 未覆盖：  DH棘轮要  求按序推进 & \textbf{完全覆盖} \\
\hline
F3 正交安全  约束 & 未覆盖 & 未覆盖 & 未覆盖 & 未覆盖 & 未覆盖 & 未覆盖 & \textbf{完全覆盖} \\
\hline
F4 过渡重叠  窗口 & 未覆盖 & 未覆盖 & 未覆盖 & 未覆盖 & 未覆盖 & 未覆盖 & \textbf{完全覆盖} \\
\hline
F5 HED  阈值分片 & 未覆盖 & 未覆盖 & 未覆盖 & 未覆盖 & 未覆盖 & 未覆盖 & \textbf{完全覆盖} \\
\hline
F6 MPTD  多路径分散 & 未覆盖 & 未覆盖 & 未覆盖 & 未覆盖 & 未覆盖 & 未覆盖 & \textbf{完全覆盖} \\
\hline
\end{longtable}

\noindent\textbf{补充对比：其他Y级文献的F1--F6覆盖分析}


\begin{longtable}{|p{2.0cm}|p{3.5cm}|p{3.5cm}|p{5.5cm}|}
\hline
\textbf{技术特征} & \textbf{S4 Hybrid KEM Combiners} & \textbf{S6 SoK: Crypto Agility} & \textbf{本申请 DMHP} \\
\hline
F1 会话内动态跳频 & 未覆盖：静态并行组合（Combiner），非序列化动态跳频 & 未覆盖：综述指出现有方案缺少会话内细粒度动态切换（Section~IV-B, pp.\,12--14） & \textbf{完全覆盖} \\
\hline
F2 确定性无状态派生 & 未覆盖 & 未覆盖 & \textbf{完全覆盖} \\
\hline
F3 正交安全约束 & 部分：多KEM组合隐含多假设安全（Section~3, Theorem~1），但无显式类别距离度量与约束机制 & 未覆盖：指出缺少可度量多样性约束（Section~VI ``Open Problems'', pp.\,21--23） & \textbf{完全覆盖} \\
\hline
F4 过渡重叠窗口 & 未覆盖 & 未覆盖 & \textbf{完全覆盖} \\
\hline
F5 HED阈值分片 & 未覆盖 & 未覆盖 & \textbf{完全覆盖} \\
\hline
F6 MPTD多路径分散 & 未覆盖 & 未覆盖 & \textbf{完全覆盖} \\
\hline
\end{longtable}

\subsection*{（四）最接近现有技术、差异特征与技术效果（审查口径化表述建议）}

\textbf{1. 最接近现有技术的确定：}
从工程实现形态看，N1（TLS 1.3及其Key Update机制）与N3（QUIC采用TLS保护并引入基于包号/密钥阶段的派生与更新）均属于"\textbf{在会话期间进行密钥更新/密钥派生}"路线，与本申请的"按序号/时间片驱动的记录/包级保护"较为接近，可作为最接近现有技术族进行对比。此外，专利P1（CN113162767A）在"算法动态切换"方面与本申请最为接近，但缺少F2--F6特征。

\textbf{2. 相对于最接近现有技术的主要差异特征：}
与上述现有技术相比，本申请至少具有如下差异特征组合（对应前述F1--F4为核心）：
\begin{itemize}[leftmargin=2em]
    \item \textbf{差异D1（对应F1）：}在同一会话内进行\textbf{跨密码原语/算法族的动态跳频}，而不仅是同一算法套件下的密钥更新；
    \item \textbf{差异D2（对应F2）：}以$\text{SeqID}$（或时间片）为输入，实现\textbf{确定性无状态派生}得到"算法索引+每单元密钥材料"，从而天然支持丢包/乱序；
    \item \textbf{差异D3（对应F3）：}引入困难问题类别（HPC）与类别距离约束$\ge d_{\min}$，实现\textbf{可度量的正交安全轮换规则}；
    \item \textbf{差异D4（对应F4）：}在切换点提供过渡重叠窗口与双解码策略，以在真实网络抖动/时延条件下保持可用性。
\end{itemize}

\textbf{3. 由差异特征产生的技术问题与技术效果：}
\begin{itemize}[leftmargin=2em]
    \item 通过D1+D3，将连续数据保护的数学假设基础进行跨类别切换，解决"单一算法/同类数学假设被削弱后导致\textbf{连续失守}"的问题，技术效果为：\textbf{降低单点突破的连续影响范围}；
    \item 通过D2，在记录/包级场景中无需依赖严格前序状态即可定位算法与密钥上下文，解决“丢包/乱序条件下会话内频繁更新导致\textbf{解密同步困难}”的问题，技术效果为：\textbf{提高乱序容忍与并行处理能力}；
    \item 通过D4，解决“切换点时钟漂移/网络抖动导致\textbf{阶段性可用性下降}”的问题，技术效果为：\textbf{在不牺牲核心安全策略的前提下提升工程可用性}；
    \item 上述效果共同作用，降低SNDL攻击的"批量采集、未来统一解密"的价值密度，技术效果为：\textbf{提高长期保密的攻击成本与不确定性}。
\end{itemize}

\textbf{4. 非显而易见性说明（组合并非简单拼接）：}
即使现有技术已公开"会话内rekey/密钥更新"（N1）以及"基于包号/阶段的密钥派生与更新"（N3），本申请的D1--D4组合仍非显而易见，原因在于：
\begin{itemize}[leftmargin=2em]
    \item 将"算法级跳频"（D1）引入记录/包级保护后，必须与无状态派生（D2）和过渡窗口（D4）共同设计，否则在乱序与重传条件下会出现解密失败率上升、重放检测与密钥更新一致性冲突等工程问题；
    \item 现有的"算法敏捷"指导（N4 RFC~7696）多为版本级/策略级治理建议，\textbf{未给出可度量的困难类别距离约束（D3）}以及与记录/包级无状态派生联动的可执行机制。事实上，RFC~7696明确倡导"版本级算法协商与迁移"（Section~3），给出的技术启示方向与本申请的"会话内逐记录跳频"恰恰相反，本领域技术人员并无动机从RFC~7696出发设计D1；
    \item S3（Signal双棘轮）有意选择固定算法以简化实现并保证前向安全性，其DH棘轮依赖严格的消息顺序交付。本领域技术人员若从S3出发，需\textbf{放弃}其核心设计假设（算法固定、顺序交付），这与S3的设计目标矛盾，不构成组合的技术启示；
    \item S6（SoK综述）在Section~VI ``Open Problems''中明确指出"\textbf{现有方案缺少可度量的密码多样性约束机制}"（pp.\,21--23），证实F3所解决的问题在现有技术中尚无可执行方案，进一步支持D3的非显而易见性；
    \item D3的"类别距离$\ge d_{\min}$"使算法多样性从"可配置"提升为"可度量、可审计、可强制执行"的规则体系，属于机制层面的实质改进，而非简单列举多算法；
    \item 最接近专利P1（CN113162767A）仅覆盖F1（粗粒度算法切换），其与F2--F4的组合需要解决无状态派生与正交约束的联动一致性问题（参见交底书第六节"技术流程四：正交安全约束与类别距离度量"中关于$A_{\text{prev}}^{\text{raw}}(s)$非递归正交校正的设计，以及"技术流程七a：可选前向安全增强（棘轮更新）"中与F2无状态派生的兼容设计），并非简单功能叠加。
\end{itemize}

\subsection*{（五）工程挑战与风险点（审查/落地的防御性说明建议）}

为提升本申请文本在审查阶段的可实施性评价，并为后续说明书实施例补充提供依据，建议在检索报告中明确以下工程挑战与对应的防御性设计要点（不改变核心创造性结论，仅用于"可实施性/可信实现"论证）：

\textbf{工程挑战与交底书技术流程/实施例的映射关系：}
\begin{center}
\begin{tabular}{|c|p{5cm}|p{6cm}|}
\hline
\textbf{编号} & \textbf{工程挑战} & \textbf{交底书对应章节} \\
\hline
E1 & 密文尺寸差异与MTU碎片化 & 技术流程三（流量整形/自适应填充） \\
\hline
E2 & MasterSecret前向安全性 & 技术流程七a（棘轮更新/迭代KDF） \\
\hline
E3 & 算法库协商与启动时延 & 系统架构(3)（MTI算法集协商） \\
\hline
E4 & 接收端状态同步与DoS风险 & 技术流程三（防重放窗口/预过滤机制） \\
\hline
E5 & 降级攻击防御（延迟包/操纵网络） & 系统架构(5)与(6)（严格序号递增与单向安全棘轮） \\
\hline
\end{tabular}
\end{center}

\begin{itemize}[leftmargin=2em]
    \item \textbf{E1 不同PQC算法的密文尺寸差异与MTU碎片化风险：}
    由于算法库可能包含格类、编码类、哈希基等不同PQC原语，其密文/封装输出长度存在显著差异。若按包/记录级跳频，输出包长可能随算法切换发生抖动，带来（i）网络层MTU碎片化、额外重传与性能下降；（ii）攻击者通过包长度统计推断当前算法类别的\textbf{侧信道}风险。
    \textbf{防御性建议：}在具体实施方式中引入\textbf{流量整形/自适应填充}模块：设定目标传输单元长度$L_\text{target}$（可取算法库中最大输出长度加冗余量），对短密文追加随机填充，使链路层观测到的包长呈现恒定或伪随机分布，从而同时减轻碎片化与"算法指纹"泄露。

    \item \textbf{E2 会话主密钥（MasterSecret）的前向安全性（PFS）与回溯风险：}
    若会话期间$\text{MasterSecret}$保持静态，则在攻击者于时刻$T$获取会话主密钥的极端场景下，可能存在对历史数据的回溯解密风险（取决于派生链路与实现细节）。
    \textbf{防御性建议：}在实现例中加入\textbf{棘轮/迭代更新}机制：每隔$N$个受保护单元或每个时间周期，执行$\text{MasterSecret}_{i+1}=\text{KDF}(\text{MasterSecret}_i,\ \text{ratchet})$并安全擦除旧值。该机制与F2无状态派生可并存：接收端在可配置窗口内保留少量近邻状态（例如最近$w$个棘轮点的种子）用于追赶与乱序容忍。\textit{注：交底书第六节"技术流程七a：可选前向安全增强（棘轮更新）"已对该机制进行详细描述，包括棘轮推进公式、与F2兼容的窗口设计及实施边界说明。}

    \item \textbf{E3 算法库协商、实现体积与启动时延：}
    若算法库覆盖多个困难问题类别，移动端/IoT设备可能面临代码体积、内存与初始化开销。
    \textbf{防御性建议：}可将算法库拆分为"必选最小集（MTI-like）+可选扩展集"，握手阶段协商可用集合ID；并允许以"类别级"协商（仅协商HPC类别与版本）降低协商长度，具体算法索引由KDF在集合内部确定。

    \item \textbf{E4 接收端状态同步与DoS攻击风险：}
    虽然F2采用无状态派生，但若攻击者注入大量无效大序号包，可能诱发接收端执行高消耗的PQC解密。
    \textbf{防御性建议：}引入\textbf{低成本预过滤机制}：在PQC解密前，基于轻量级对称MAC或Bloom Filter对$\text{SeqID}$的有效性及窗口范围进行预校验，快速丢弃非法流量，保护接收端计算资源。

    \item \textbf{E5 降级攻击防御（延迟包/操纵网络）：}
    攻击者可能通过故意延迟数据包来延长重叠窗口，或通过注入网络延迟/丢包来操纵威胁指标，迫使系统降级到较弱的算法或状态。
    \textbf{防御性建议：}在重叠窗口机制中引入\textbf{严格序号递增与即时关闭策略}：一旦成功验证并处理了新算法的数据包，立即关闭旧算法的接收窗口。对于威胁感知模块，引入\textbf{单向安全棘轮}：系统可自动升级到高安全模式，但降级必须经过持续的干净网络状态验证或显式的密码认证控制消息，且算法选择的确定性和正交性不受跳频频率改变的影响。
\end{itemize}

        \textbf{1. 与"算法协商与密码敏捷"相关公开资料的对比：}
传统协议通常在握手阶段确定并锁定单一（或少量）算法套件，在会话期间长时间复用。其弱点在于算法与实现风险集中、容易被SNDL批量采集并等待未来突破。本发明的DMHP在会话内细粒度跳频，降低任一算法长期暴露的价值，且可在不重建会话的前提下完成动态切换。

        \textbf{2. 与“后量子标准化与单算法迁移”相关公开资料的对比：}
单一PQC算法替换可提升抗量子能力，但无法消除算法族在未来被削弱的不确定性。本发明通过“硬问题类别正交约束”，在结构化格、编码理论、哈希基等类别间切换，形成可度量的多样性防线。

        \textbf{3. 与SNDL威胁分析相关公开资料的对比：}
SNDL核心在于长期保存密文等待未来大规模算力/新攻击。本发明通过微碎片化（按包/按块/按分片）减少单一算法保护的连续数据规模，并可选通过阈值分片（HED）将重构门槛提升为“至少解出$k$个份额”。

        \textbf{4. 与多路径传输协议相关公开资料的对比：}
多路径协议通常解决吞吐/可靠性或抗链路故障问题。本发明将多路径与算法跳频结合，形成“空间正交”：攻击者需要同时捕获多条物理路径并跨多个算法类别破译，显著提高完整会话重构难度。

% 五、检索结论
\section*{五、检索结论}

基于初始检索（2026.01.26）及补充更新检索（2026.02.12）对包括RFC标准、IETF草案、专利文献及经典扩频理论等文献的全面检索与分析，结论如下：

\textbf{未检索到X级（特别相关）文献}：在所有检索范围内，未发现单一公开文本同时覆盖本申请的"会话内动态跳频（F1）"、"确定性无状态派生（F2）"、"正交安全约束（F3）"与"过渡重叠窗口（F4）"等核心特征，即不存在可单独影响本申请新颖性判断的X级文献。

现有技术主要存在以下局限：

\begin{itemize}[leftmargin=2em]
    \item 握手阶段算法协商/迁移（密码敏捷），但缺少会话内细粒度跳频与工程可用的过渡机制（F1/F4）；
    \item 基于单一PQC算法族的替换或混合，但缺少跨困难类别的“正交安全约束”与可度量的多样性策略（F3）；
    \item 多路径传输以性能/可靠性为主，缺少与算法跳频及阈值分片联动的安全分散框架（F6与F1/F5联动）。
\end{itemize}

\textbf{综上所述}，在所有检索范围内，未发现单一公开文献同时覆盖本申请的F1--F4核心特征（即不存在X级文献），\textbf{本申请核心技术方案具备新颖性}。其F1--F4的有机组合相对于已检索到的Y级文献的任意两两组合，均需要解决无状态派生与正交约束的联动一致性、乱序条件下的解密同步、以及切换点的工程可用性等非显而易见的技术问题，\textbf{该组合具备非显而易见性（创造性）}。具体而言：

\begin{itemize}[leftmargin=2em]
    \item \textbf{F2（无状态派生/乱序容忍）与F4（过渡重叠窗口/双解码）并非简单拼接：}其需要在允许乱序和重传的网络条件下同时满足解密成功率、重放检测与密钥更新一致性，属于记录层/包级密钥更新的工程约束体系；
    \item \textbf{F3（困难问题类别距离约束）区别于一般"支持多算法/可配置"的密码敏捷：}本申请提供可度量的"类别距离$\ge d_{\min}$"选择规则（参见交底书中的距离矩阵配置，其中$d(\text{Lattice-S}, \text{Lattice-U})=0.5$, $d(\text{Code}, \text{Hash})=0.8$, 跨大类均为$1.0$），使连续受保护单元跨数学假设类别进行轮换，从而降低同类突破造成连续失守的风险。
\end{itemize}

建议：后续审查中可围绕"会话内频繁rekey/记录层动态切换/阈值分片重构/多路径安全传输结合"等主题进一步扩展检索，以排除相关性较低的组合干扰。

\end{document}
